\section*{Problema 2}

Encuentra la ecuación de la recta en el espacio 2D, perpendicular a $\vec{A}$ y que pasa por $\vec{P}$, para los siguientes valores de $\vec{A}$ y $\vec{P}$:

3. $\vec{A} = (1, -1)$, $\vec{P} = (-5, 3)$

\textbf{Solución:}

La ecuación de una recta perpendicular a un vector $\vec{A} = (a, b)$ y que pasa por un punto $\vec{P} = (x_0, y_0)$ se puede escribir como:

\[
ax + by = c
\]

donde $c$ se determina sustituyendo las coordenadas del punto $\vec{P}$ en la ecuación. En este caso, $\vec{A} = (1, -1)$, por lo que la ecuación de la recta perpendicular es:

\[
1 \cdot x - 1 \cdot y = c
\]

Sustituyendo $\vec{P} = (-5, 3)$:

\[
1(-5) - 1(3) = c \implies c = -8
\]

Por lo tanto, la ecuación de la recta es:

\[
x - y = -8
\]

4. $\vec{A} = (-5, 4)$, $\vec{P} = (3, 2)$

\textbf{Solución:}

De manera similar, para $\vec{A} = (-5, 4)$, la ecuación de la recta perpendicular es:

\[
-5x + 4y = c
\]

Sustituyendo $\vec{P} = (3, 2)$:

\[
-5(3) + 4(2) = c \implies c = -7
\]

Por lo tanto, la ecuación de la recta es:

\[
-5x + 4y = -7
\]
